% ============================================================
% Paper_GFJ — When Geopolitical Shocks Become Systemic
% Target: Global Finance Journal
% ============================================================
\documentclass[11pt]{article}

\usepackage[T1]{fontenc}
\usepackage[utf8]{inputenc}
\usepackage{lmodern}
\usepackage{microtype}
\usepackage[margin=1in]{geometry}

\usepackage{graphicx}
\graphicspath{
  {figures/}
  {../results/figures/}
  {../results/fragility/}
  {../results/connectedness/}
  {../results/emfi/}
  {../results/stress_regimes/}
  {../results/panel_lp/}
  {../results/state_lp/}
  {../results/event_classification/}
}

\usepackage{booktabs}
\usepackage{threeparttable}
\usepackage{array}
\usepackage{longtable}

\usepackage{amsmath,amssymb}

\usepackage[strings]{underscore}

\usepackage[hidelinks]{hyperref}
\usepackage[nameinlink,capitalise]{cleveref}

\usepackage[numbers,sort&compress]{natbib}

% ---- Working title; update when finalised ----
\title{When Geopolitical Shocks Become Systemic:\\
  Volatility Connectedness and Market-Implied Stress\\
  in European Equity Markets}
\author{[Author information omitted for blind review]}
\date{\today}

% ============================================================
\begin{document}
\maketitle

% ---- Abstract -----------------------------------------------
\begin{abstract}
% TODO: write after main results are known.
%
% Template:
% This paper studies whether geopolitical shocks become systemic
% financial stress events in European equity markets.
% Using daily data for 19 equity markets from 2017--2025, we combine
% a validated high-frequency geopolitical conflict shock series with
% market-based measures of financial fragility, volatility
% connectedness, and market-implied stress regimes.
% [Main results to be inserted here.]
% We find that geopolitical shocks increase the European Market
% Fragility Index primarily when they amplify volatility connectedness
% or coincide with already-fragile market conditions.
\end{abstract}

\newpage
\tableofcontents
\newpage

% ============================================================
\section{Introduction}
\label{sec:introduction}

% TODO: write after empirical results are finalised.
%
% Motivate: geopolitical risk -- financial stability link
% Gap: existing literature studies return/volatility, not systemic
%   fragility transmission and connectedness
% This paper: three-layer approach
%   1. Validated shock trigger (EVT+FDR series from Lupu et al. 2026)
%   2. Daily systemic fragility metrics + EMFI
%   3. HMM market-implied stress classification
% Contributions (three):
%   (i)   novel EMFI daily index for European equity markets
%   (ii)  evidence that shocks amplify connectedness, not just volatility
%   (iii) typology of absorbed vs. systemic geopolitical shocks

% ============================================================
\section{Related Literature}
\label{sec:literature}

% TODO

% ============================================================
\section{Data and Geopolitical Shock Triggers}
\label{sec:data}

\subsection{Equity return panel}
% 19 European equity markets, 2017-01-01 to 2025-10-15.
% Daily log returns from LSEG Workspace.
% Table~\ref{tab:data_summary}: country coverage, index names, cluster assignment.

\subsection{Geopolitical conflict shock series}
% Reuse EVT+FDR shock series from companion paper (Lupu et al. 2026, BIR).
% Briefly describe: GDELT Doc 2.0, 28d WMA, GPD tail, FDR cross-section, declustering.
% Key statistic: N detected shock events, range of S values.
% Aggregate shock: MaxShock_t = max_i S_{i,t}.

% ============================================================
\section{Market-Based Financial Stability Measures}
\label{sec:fragility}

\subsection{Daily volatility stress and tail co-exceedance}
% VolStress_t, TailVolBreadth_t, TailLossBreadth_t, AvgCorr_t
% Rolling threshold construction: no look-ahead, 252-day backward window.

\subsection{Volatility connectedness index}
% Rolling Diebold-Yilmaz VAR-GFEVD on |r_{i,t}|.
% Window W=100 days; VAR order by BIC capped at p=3; H=10 step horizon.
% Report TCI_t, plus directional spillovers.

\subsection{European Market Fragility Index (EMFI)}
% PCA over [VolStress, TailVolBreadth, AvgCorr, TCI].
% First PC; orientation: higher = more fragile.
% Report: loadings, variance explained.
% Figure~\ref{fig:emfi_timeline}: EMFI over time with event overlays.

% ============================================================
\section{Market-Implied Systemic Stress Regimes}
\label{sec:hmm}

% Gaussian HMM, 3 states, estimated on X_t = [VolStress, TailVolBreadth, AvgCorr, TCI].
% Key: estimated without reference to geopolitical shock dates.
% P_stress_t = daily probability of systemic-stress state.
% Validation against known events.
% Figure~\ref{fig:hmm_regimes}.

% ============================================================
\section{Dynamic Effects of Geopolitical Shocks on Systemic Fragility}
\label{sec:lp}

\subsection{Empirical design}
% Local projections (Jord\`a 2005).
% Specification:
%   Y_{t+k} = alpha + beta_k MaxShock_t + gamma Y_{t-1} + delta r^{global}_t + mu_m + eps
% Horizons k=-5 to +15. Newey-West SEs.

\subsection{Baseline results}
% Table + Figure for each outcome: EMFI, TCI, TailBreadth, P_stress, AvgCorr.
% Figure~\ref{fig:lp_combined}.

% ============================================================
\section{State Dependence and Event Classification}
\label{sec:state}

\subsection{State-dependent local projections}
% HighFragility_{t-1} indicator (EMFI > 75th pctile).
% Interaction specification.
% Figure~\ref{fig:state_lp}: normal vs. fragile IRF comparison.

\subsection{Geopolitical event taxonomy}
% Four categories: Absorbed, Localized, Systemic, Market-only.
% Table with major events classified.
% Figure~\ref{fig:network_events}: connectedness network around top Systemic events.

% ============================================================
\section{Robustness}
\label{sec:robustness}

% Checks: alternative connectedness, alternative shock aggregation,
% alternative EMFI weights, COVID/Ukraine exclusions, future-shock placebo.

% ============================================================
\section{Conclusion}
\label{sec:conclusion}

% TODO

% ============================================================
\bibliographystyle{plainnat}
\bibliography{references}

% ============================================================
\appendix

\section{Data Details}
\label{app:data}

% Table A1: Data summary (countries, date ranges, index names)

\section{EMFI Construction Details}
\label{app:emfi}

% PCA loadings table, scree plot, robustness (equal-weight alternative)

\section{HMM Parameter Estimates}
\label{app:hmm}

% Transition matrix, emission parameters per state

\section{Additional LP Results}
\label{app:lp}

% Individual outcome LP tables (numerical coefficients + SEs)

\section{Robustness Tables}
\label{app:robustness}

\end{document}
